\chapter{Cäsar-Chiffre (Python)}
Das \textbf{Caesar-Chiffre} ist eine einfache Verschlüsselung, bei der jeder Buchstabe im Text um eine feste Zahl von Stellen im Alphabet verschoben wird.\par
Bei einer Verschiebung von \ac{zB} 3 wird aus „A“ „D“, aus „B“ „E“ und so weiter. Um den Text zu entschlüsseln, wird der Schlüssel (die Zahl der Verschiebung) einfach in die entgegengesetzte Richtung angewendet.
Es ist einfach, aber unsicher, da es nur 25 mögliche Schlüssel gibt.

\subsection*{Cäsar-Chiffre - A}
Teil A verschlüsselt oder entschlüsselt Text mit dem Cäsar-Verschlüsselungsverfahren, bei dem Buchstaben im Alphabet um eine bestimmte Anzahl (den „Schlüssel“) verschoben werden. Es funktioniert über die Kommandozeile.\\\\
\textbf{Eingabe:} python CaesarCipher-1-a.py "text" key
\begin{codeblock}
    \lstinputlisting[
        style=inputlistingstyle,
        caption={Cäsar-Chiffre - A},
        label={lst:CaesarCipher-1-a}
    ]{code/CaesarCipher-1-a.py}
\end{codeblock}

\subsection*{Cäsar-Chiffre - B}
In diesem Teil wird gezählt, wie oft jeder Buchstabe in einem gegebenen Text vorkommt (Groß- und Kleinschreibung werden dabei ignoriert). Das Ergebnis ist ein Histogramm in Form eines Wörterbuchs.\\\\
\textbf{Ausgabe von Text:}\quad \{'t': 2, 'e': 1, 's': 1\}
\begin{codeblock}
    \lstinputlisting[
        style=inputlistingstyle,
        caption={Cäsar-Chiffre - B},
        label={lst:CaesarCipher-1-b}
    ]{code/CaesarCipher-1-b.py}
\end{codeblock}

\subsection*{Cäsar-Chiffre - C}
Dieses Programm berechnet die Häufigkeit von Buchstaben in einem gegebenen Text und berechnet dann die Wahrscheinlichkeiten für jeden Buchstaben (a-z). Die Ergebnisse werden als Prozentsatz der gesamten Buchstabenanzahl ausgegeben.\\\\
\textbf{Ausgabe:}\quad Buchstabenhäufigkeit (in \%):\\
a: 0.00\%\\
b: 0.00\%\\
c: 0.00\%\\
d: 0.00\%\\
e: 25.00\%\\
f: 0.00\%\\
g: 0.00\%\\
h: 0.00\%\\
i: 0.00\%\\
j: 0.00\%\\
k: 0.00\%\\
l: 0.00\%\\
m: 0.00\%\\
n: 0.00\%\\
o: 0.00\%\\
p: 0.00\%\\
q: 0.00\%\\
r: 0.00\%\\
s: 25.00\%\\
t: 50.00\%\\
u: 0.00\%\\
v: 0.00\%\\
w: 0.00\%\\
x: 0.00\%\\
y: 0.00\%\\
z: 0.00\%\\
\begin{codeblock}
    \lstinputlisting[
        style=inputlistingstyle,
        caption={Cäsar-Chiffre - C},
        label={lst:CaesarCipher-1-c}
    ]{code/CaesarCipher-1-c.py}
\end{codeblock}

\subsection*{Cäsar-Chiffre - D}
Im letzten Teil nutzt man den \textbf{Cäsar-Chiffre-Algorithmus} zur Verschlüsselung und Entschlüsselung von Texten. Es analysiert auch die \textbf{Häufigkeit} von Buchstaben im Text, um die \textbf{Wahrscheinlichkeiten} der Buchstaben zu berechnen. Zudem enthält es eine Methode, um den verschlüsselten Text mit \textbf{Chi-Quadrat-Test} zu entschlüsseln.\\\\
\textbf{Ausgabe von Test:}\\
Verschlüsselter Text:\\
Whvw\par
Häufigkeit vorkommender Buchstaben:\\
\{'t': 2, 'e': 1, 's': 1\}\par
Wahrscheinlichkeiten der Buchstaben im Text:\\
e=25.0\%, s=25.0\%, t=50.0\%, d-x=0\%\par
Entschlüsselter Text:\\
Test\\
\begin{codeblock}
    \lstinputlisting[
        style=inputlistingstyle,
        caption={Cäsar-Chiffre - D},
        label={lst:CaesarCipher-1-d}
    ]{code/CaesarCipher-1-d.py}
\end{codeblock}

\subsection*{Zusammenfassung A-C}
\begin{itemize}[itemsep=0pt]
   \item \textbf{Teil A (Cäsar-Chiffre)} 
   \begin{itemize}
    \item Initialisiert ein Array mit 8 leeren BitMessage-Objekten (für 8 Bits)...8 BitMessage-Instanzen = ein Zeichens
   \end{itemize}
   \item \textbf{Teil B (Häufigkeitsanalyse)} 
   \begin{itemize}
    \item Zählt, wie oft jeder Buchstabe im Text vorkommt, und berechnet die Wahrscheinlichkeit jedes Buchstabens
   \end{itemize}
   \item \textbf{Teil C (Chi-Quadrat-Analyse)} 
   \begin{itemize}
    \item Entschlüsselt einen Text basierend auf der Häufigkeit der Buchstaben und vergleicht diese mit einem Beispieltext, um den besten Schlüssel zu finden.
   \end{itemize}
\end{itemize}